\documentclass[11pt]{article}
\usepackage[margin=1in]{geometry}
\usepackage{amsmath,amssymb}
\usepackage{microtype}
\usepackage{xcolor}
\usepackage[colorlinks=true,linkcolor=blue!55!black,citecolor=blue!55!black,urlcolor=blue!55!black]{hyperref}

\title{The LUR quasihyperbolic-ball question in arXiv:1407.2403\\
\large Explicitly answered by Theorem 3.2 of arXiv:1708.02240}
\author{Automated Mathematics Run \texttt{fa\_banach\_001}, lane 2}
\date{19 July 2026}

\begin{document}
\maketitle

\begin{center}
\fcolorbox{blue!55!black}{blue!4}{\parbox{0.88\textwidth}{
\textbf{Status.} \texttt{literature\_already\_answered}.  This is an exact,
author-acknowledged later answer, not a new solution from the run.}}
\end{center}

\section*{Original question}

Kl\'en, Rasila, and Talponen ask on source PDF page 15, at the end of
Section 4 of \emph{On smoothness of quasihyperbolic balls}~\cite{KRT}, whether
centered quasihyperbolic balls in symmetric convex domains of reflexive LUR
Banach spaces induce LUR norms through their Minkowski functionals.  In the
notation of the paper, for $B=B_k(0,r)$ this functional is
\[
 M_B(x)=\inf\{\lambda>0:x\in\lambda B\}.
\]

\section*{Exact later answer}

Rasila, Talponen, and Zhang state in their abstract that they answer a question
posed by the first two authors with R. Kl\'en~\cite{RTZ}.  Immediately before
Theorem 3.2 on supporting PDF page 5, they identify the LUR Minkowski question
from~\cite{KRT} and say it is settled affirmatively in the reflexive case.
Theorem 3.2 proves the stronger sequential statement: if $X$ is reflexive and
LUR, $\Omega\subset X$ is convex, and
\[
 k(x_0,y)=k(x_0,y_n),\qquad
 k\!\left(x_0,\frac{y+y_n}{2}\right)\longrightarrow k(x_0,y),
\]
then $y_n\to y$ in norm.  Its final sentence gives the requested conclusion
verbatim in mathematical content: if $\Omega$ is symmetric, the Minkowski
functional of $B_k(0,r)$ is an equivalent LUR norm for every $r>0$.

Thus the question extracted from arXiv:1407.2403 has a complete positive
literature answer in Theorem 3.2 of arXiv:1708.02240.  The overlap is explicit,
not an agent-inferred reformulation.

\section*{Search bounds and scope}

The run's lightweight indexes had no packet for this exact question.  A bounded
search through 19 July 2026 used arXiv:1407.2403, the exact question phrase,
``quasihyperbolic balls LUR Minkowski functional'', and ``reflexive LUR
quasihyperbolic''.  It found arXiv:1708.02240 and no need for a new proof.
The supporting theorem exactly covers the original reflexive-LUR question;
no part of that stated question remains open in this packet.

\begin{thebibliography}{9}
\bibitem{KRT}
R. Kl\'en, A. Rasila, and J. Talponen,
\emph{On smoothness of quasihyperbolic balls},
Annales Academiae Scientiarum Fennicae Mathematica 42 (2017), 439--452;
arXiv:1407.2403.

\bibitem{RTZ}
A. Rasila, J. Talponen, and X. Zhang,
\emph{Observations on quasihyperbolic geometry modeled on Banach spaces},
arXiv:1708.02240 (2017), Theorem 3.2.
\end{thebibliography}

\end{document}
